%%%%%%%%%%%%%%%%%%%%%%%%%%%%%%%%%%%%%%%%%%%%%%%%%%%%%%%%%%%%
% 6.4 GENERATING FUNCTIONS
% The Composition of Advanced Mathematics
%%%%%%%%%%%%%%%%%%%%%%%%%%%%%%%%%%%%%%%%%%%%%%%%%%%%%%%%%%%%

\section{Generating Functions}
\label{sec:generating-functions}

\prerequisite{
Sequences, power series notation, algebraic manipulation, recurrence
relations, binomial coefficients, and basic counting.
}

\learningobjectives{
By the end of this section, the reader should be able to:
\begin{itemize}
    \item define ordinary and exponential generating functions;
    \item translate sequences into formal power series;
    \item recover coefficients from generating functions;
    \item perform algebraic operations on generating functions;
    \item interpret products using convolution;
    \item use geometric-series identities;
    \item derive generating functions from recurrence relations;
    \item solve linear recurrences using generating functions;
    \item derive generating functions for Fibonacci-type sequences;
    \item use differentiation and integration to transform sequences;
    \item apply index shifts correctly;
    \item use generating functions to count restricted selections;
    \item interpret products combinatorially;
    \item use generating functions for integer partitions;
    \item distinguish ordinary from exponential generating functions;
    \item recognize when exponential generating functions are natural;
    \item use partial fractions to recover explicit coefficient formulas;
    \item understand multivariate generating functions;
    \item recognize generating functions as algebraic encodings of
          combinatorial classes.
\end{itemize}
}

%%%%%%%%%%%%%%%%%%%%%%%%%%%%%%%%%%%%%%%%%%%%%%%%%%%%%%%%%%%%
\subsection{Why Generating Functions?}
\label{subsec:why-generating-functions}
%%%%%%%%%%%%%%%%%%%%%%%%%%%%%%%%%%%%%%%%%%%%%%%%%%%%%%%%%%%%

A sequence

\[
a_0,a_1,a_2,\ldots
\]

contains infinitely many values.

A generating function packages the entire sequence into a single formal
power series.

The basic idea is

\[
\boxed{
a_0,a_1,a_2,\ldots
\quad\longleftrightarrow\quad
a_0+a_1x+a_2x^2+\cdots.
}
\]

The variable \(x\) acts primarily as a bookkeeping device.

Once the sequence is encoded algebraically, operations on power series
can reveal:

\begin{itemize}
    \item recurrence relations;
    \item explicit formulas;
    \item convolution identities;
    \item counting formulas;
    \item asymptotic information;
    \item structural decompositions.
\end{itemize}

%%%%%%%%%%%%%%%%%%%%%%%%%%%%%%%%%%%%%%%%%%%%%%%%%%%%%%%%%%%%
\subsection{Ordinary Generating Functions}
\label{subsec:ordinary-generating-functions}
%%%%%%%%%%%%%%%%%%%%%%%%%%%%%%%%%%%%%%%%%%%%%%%%%%%%%%%%%%%%

\begin{definitionbox}{Ordinary Generating Function}
For a sequence

\[
(a_n)_{n\geq0},
\]

its ordinary generating function, abbreviated OGF, is

\[
\boxed{
A(x)
=
\sum_{n=0}^{\infty}a_nx^n.
}
\]
\end{definitionbox}

Thus,

\[
A(x)
=
a_0+a_1x+a_2x^2+a_3x^3+\cdots.
\]

The coefficient of \(x^n\) records the value \(a_n\).

%%%%%%%%%%%%%%%%%%%%%%%%%%%%%%%%%%%%%%%%%%%%%%%%%%%%%%%%%%%%
\subsection{Coefficient Extraction}
\label{subsec:coefficient-extraction}
%%%%%%%%%%%%%%%%%%%%%%%%%%%%%%%%%%%%%%%%%%%%%%%%%%%%%%%%%%%%

If

\[
A(x)=\sum_{n\geq0}a_nx^n,
\]

then the notation

\[
\boxed{
[x^n]A(x)=a_n
}
\]

means ``the coefficient of \(x^n\) in \(A(x)\).''

For example, if

\[
A(x)=1+4x+7x^2+10x^3+\cdots,
\]

then

\[
[x^2]A(x)=7.
\]

Coefficient-extraction notation is extremely useful in advanced
combinatorics.

%%%%%%%%%%%%%%%%%%%%%%%%%%%%%%%%%%%%%%%%%%%%%%%%%%%%%%%%%%%%
\subsection{Generating Function of a Constant Sequence}
\label{subsec:gf-constant-sequence}
%%%%%%%%%%%%%%%%%%%%%%%%%%%%%%%%%%%%%%%%%%%%%%%%%%%%%%%%%%%%

Consider the sequence

\[
1,1,1,1,\ldots.
\]

Its generating function is

\[
1+x+x^2+x^3+\cdots.
\]

Using the geometric-series identity,

\[
\boxed{
\frac{1}{1-x}
=
\sum_{n=0}^{\infty}x^n.
}
\]

Therefore,

\[
\boxed{
[x^n]\frac{1}{1-x}=1.
}
\]

%%%%%%%%%%%%%%%%%%%%%%%%%%%%%%%%%%%%%%%%%%%%%%%%%%%%%%%%%%%%
\subsection{Formal Power Series}
\label{subsec:formal-power-series}
%%%%%%%%%%%%%%%%%%%%%%%%%%%%%%%%%%%%%%%%%%%%%%%%%%%%%%%%%%%%

In combinatorics, generating functions are often treated as
\emph{formal power series}.

This means that

\[
A(x)
=
\sum_{n\geq0}a_nx^n
\]

is manipulated algebraically without requiring numerical convergence.

The symbol \(x\) serves as a formal marker.

Thus identities such as

\[
\frac{1}{1-x}
=
1+x+x^2+x^3+\cdots
\]

may be interpreted coefficientwise.

\begin{importantbox}
For formal generating functions, the primary question is not whether
the infinite series converges numerically.

The primary question is whether the coefficients are manipulated
correctly.
\end{importantbox}

%%%%%%%%%%%%%%%%%%%%%%%%%%%%%%%%%%%%%%%%%%%%%%%%%%%%%%%%%%%%
\subsection{Basic Generating Function Dictionary}
\label{subsec:gf-dictionary}
%%%%%%%%%%%%%%%%%%%%%%%%%%%%%%%%%%%%%%%%%%%%%%%%%%%%%%%%%%%%

Several standard sequences have simple generating functions.

\[
\boxed{
\begin{array}{c|c}
\text{Sequence }a_n & \text{Generating Function}\\
\hline
1 & \dfrac{1}{1-x}\\[3mm]
r^n & \dfrac{1}{1-rx}\\[3mm]
n & \dfrac{x}{(1-x)^2}\\[3mm]
n+1 & \dfrac{1}{(1-x)^2}\\[3mm]
\binom{n+k-1}{k-1}
&
\dfrac{1}{(1-x)^k}
\end{array}
}
\]

These identities form a basic computational toolkit.

%%%%%%%%%%%%%%%%%%%%%%%%%%%%%%%%%%%%%%%%%%%%%%%%%%%%%%%%%%%%
\subsection{Geometric Generating Functions}
\label{subsec:geometric-generating-functions}
%%%%%%%%%%%%%%%%%%%%%%%%%%%%%%%%%%%%%%%%%%%%%%%%%%%%%%%%%%%%

For the sequence

\[
a_n=r^n,
\]

we have

\[
A(x)
=
\sum_{n=0}^{\infty}r^nx^n.
\]

Since

\[
r^nx^n=(rx)^n,
\]

this becomes

\[
\boxed{
A(x)
=
\frac{1}{1-rx}.
}
\]

Therefore,

\[
\boxed{
[x^n]\frac{1}{1-rx}=r^n.
}
\]

%%%%%%%%%%%%%%%%%%%%%%%%%%%%%%%%%%%%%%%%%%%%%%%%%%%%%%%%%%%%
\subsection{Generating Functions and Shifts}
\label{subsec:gf-shifts}
%%%%%%%%%%%%%%%%%%%%%%%%%%%%%%%%%%%%%%%%%%%%%%%%%%%%%%%%%%%%

Suppose

\[
A(x)=\sum_{n=0}^{\infty}a_nx^n.
\]

Multiplying by \(x\) gives

\[
xA(x)
=
\sum_{n=0}^{\infty}a_nx^{n+1}.
\]

Equivalently,

\[
xA(x)
=
\sum_{n=1}^{\infty}a_{n-1}x^n.
\]

Thus multiplication by \(x\) shifts coefficients one position to the
right.

More generally,

\[
\boxed{
x^kA(x)
=
\sum_{n=k}^{\infty}a_{n-k}x^n.
}
\]

%%%%%%%%%%%%%%%%%%%%%%%%%%%%%%%%%%%%%%%%%%%%%%%%%%%%%%%%%%%%
\subsection{Addition of Generating Functions}
\label{subsec:addition-generating-functions}
%%%%%%%%%%%%%%%%%%%%%%%%%%%%%%%%%%%%%%%%%%%%%%%%%%%%%%%%%%%%

Suppose

\[
A(x)
=
\sum_{n\geq0}a_nx^n
\]

and

\[
B(x)
=
\sum_{n\geq0}b_nx^n.
\]

Then

\[
A(x)+B(x)
=
\sum_{n\geq0}(a_n+b_n)x^n.
\]

Therefore,

\[
\boxed{
[x^n](A(x)+B(x))
=
a_n+b_n.
}
\]

Addition acts coefficientwise.

%%%%%%%%%%%%%%%%%%%%%%%%%%%%%%%%%%%%%%%%%%%%%%%%%%%%%%%%%%%%
\subsection{Scalar Multiplication}
\label{subsec:gf-scalar-multiplication}
%%%%%%%%%%%%%%%%%%%%%%%%%%%%%%%%%%%%%%%%%%%%%%%%%%%%%%%%%%%%

For a constant \(c\),

\[
cA(x)
=
\sum_{n\geq0}ca_nx^n.
\]

Hence,

\[
\boxed{
[x^n](cA(x))=ca_n.
}
\]

Linear combinations of sequences correspond directly to linear
combinations of generating functions.

%%%%%%%%%%%%%%%%%%%%%%%%%%%%%%%%%%%%%%%%%%%%%%%%%%%%%%%%%%%%
\subsection{Products and Convolution}
\label{subsec:gf-products-convolution}
%%%%%%%%%%%%%%%%%%%%%%%%%%%%%%%%%%%%%%%%%%%%%%%%%%%%%%%%%%%%

The product of generating functions is especially important.

Let

\[
A(x)=\sum_{n\geq0}a_nx^n
\]

and

\[
B(x)=\sum_{n\geq0}b_nx^n.
\]

Then

\[
A(x)B(x)
=
\sum_{n\geq0}
\left(
\sum_{k=0}^{n}a_kb_{n-k}
\right)x^n.
\]

Thus,

\[
\boxed{
[x^n](A(x)B(x))
=
\sum_{k=0}^{n}a_kb_{n-k}.
}
\]

The inner sum is called a convolution.

%%%%%%%%%%%%%%%%%%%%%%%%%%%%%%%%%%%%%%%%%%%%%%%%%%%%%%%%%%%%
\subsection{Combinatorial Meaning of Convolution}
\label{subsec:combinatorial-convolution}
%%%%%%%%%%%%%%%%%%%%%%%%%%%%%%%%%%%%%%%%%%%%%%%%%%%%%%%%%%%%

Suppose \(a_k\) counts structures of size \(k\) from one class and
\(b_j\) counts structures of size \(j\) from another.

If an object of total size \(n\) is formed by choosing:

\[
\text{one object of size }k
\]

and

\[
\text{one object of size }n-k,
\]

then the number of pairs is

\[
\sum_{k=0}^{n}a_kb_{n-k}.
\]

Thus multiplication of generating functions corresponds naturally to
combining independent structures whose sizes add.

%%%%%%%%%%%%%%%%%%%%%%%%%%%%%%%%%%%%%%%%%%%%%%%%%%%%%%%%%%%%
\subsection{Worked Example: Squaring the Geometric Generating Function}
\label{subsec:worked-geometric-square}
%%%%%%%%%%%%%%%%%%%%%%%%%%%%%%%%%%%%%%%%%%%%%%%%%%%%%%%%%%%%

\begin{workedexamplebox}{Finding Coefficients of \((1-x)^{-2}\)}

We know

\[
\frac{1}{1-x}
=
\sum_{n\geq0}x^n.
\]

Therefore,

\[
\frac{1}{(1-x)^2}
=
\left(
\sum_{n\geq0}x^n
\right)^2.
\]

The coefficient of \(x^n\) is

\[
\sum_{k=0}^{n}1.
\]

There are \(n+1\) terms in this sum.

\begin{answerbox}
\[
\boxed{
\frac{1}{(1-x)^2}
=
\sum_{n=0}^{\infty}(n+1)x^n.
}
\]
\end{answerbox}

\end{workedexamplebox}

%%%%%%%%%%%%%%%%%%%%%%%%%%%%%%%%%%%%%%%%%%%%%%%%%%%%%%%%%%%%
\subsection{The Negative Binomial Series}
\label{subsec:negative-binomial-series}
%%%%%%%%%%%%%%%%%%%%%%%%%%%%%%%%%%%%%%%%%%%%%%%%%%%%%%%%%%%%

More generally,

\[
\boxed{
\frac{1}{(1-x)^k}
=
\sum_{n=0}^{\infty}
\binom{n+k-1}{k-1}x^n.
}
\]

Equivalently,

\[
\boxed{
[x^n](1-x)^{-k}
=
\binom{n+k-1}{k-1}.
}
\]

This identity connects generating functions directly with stars and
bars.

%%%%%%%%%%%%%%%%%%%%%%%%%%%%%%%%%%%%%%%%%%%%%%%%%%%%%%%%%%%%
\subsection{Differentiating Generating Functions}
\label{subsec:differentiating-generating-functions}
%%%%%%%%%%%%%%%%%%%%%%%%%%%%%%%%%%%%%%%%%%%%%%%%%%%%%%%%%%%%

If

\[
A(x)
=
\sum_{n=0}^{\infty}a_nx^n,
\]

then formally,

\[
A'(x)
=
\sum_{n=1}^{\infty}na_nx^{n-1}.
\]

Multiplying by \(x\),

\[
\boxed{
xA'(x)
=
\sum_{n=0}^{\infty}na_nx^n.
}
\]

Thus differentiation introduces factors of \(n\).

%%%%%%%%%%%%%%%%%%%%%%%%%%%%%%%%%%%%%%%%%%%%%%%%%%%%%%%%%%%%
\subsection{Generating Function for \(n\)}
\label{subsec:gf-for-n}
%%%%%%%%%%%%%%%%%%%%%%%%%%%%%%%%%%%%%%%%%%%%%%%%%%%%%%%%%%%%

Begin with

\[
\frac{1}{1-x}
=
\sum_{n=0}^{\infty}x^n.
\]

Differentiate:

\[
\frac{1}{(1-x)^2}
=
\sum_{n=1}^{\infty}nx^{n-1}.
\]

Multiply by \(x\):

\[
\boxed{
\frac{x}{(1-x)^2}
=
\sum_{n=0}^{\infty}nx^n.
}
\]

Hence,

\[
\boxed{
[x^n]\frac{x}{(1-x)^2}=n.
}
\]

%%%%%%%%%%%%%%%%%%%%%%%%%%%%%%%%%%%%%%%%%%%%%%%%%%%%%%%%%%%%
\subsection{Generating Function for \(n^2\)}
\label{subsec:gf-for-n-squared}
%%%%%%%%%%%%%%%%%%%%%%%%%%%%%%%%%%%%%%%%%%%%%%%%%%%%%%%%%%%%

Starting from

\[
\sum_{n\geq0}nx^n
=
\frac{x}{(1-x)^2},
\]

differentiate and multiply by \(x\).

This gives

\[
\sum_{n\geq0}n^2x^n
=
\boxed{
\frac{x(1+x)}{(1-x)^3}.
}
\]

Thus polynomial sequences correspond to rational generating functions
whose denominators are powers of \(1-x\).

%%%%%%%%%%%%%%%%%%%%%%%%%%%%%%%%%%%%%%%%%%%%%%%%%%%%%%%%%%%%
\subsection{Integration of Generating Functions}
\label{subsec:integration-generating-functions}
%%%%%%%%%%%%%%%%%%%%%%%%%%%%%%%%%%%%%%%%%%%%%%%%%%%%%%%%%%%%

Formal integration gives

\[
\int A(x)\,dx
=
\sum_{n\geq0}
\frac{a_n}{n+1}x^{n+1}.
\]

Thus integration introduces division by the index.

Differentiation and integration provide systematic ways to transform
coefficient sequences.

%%%%%%%%%%%%%%%%%%%%%%%%%%%%%%%%%%%%%%%%%%%%%%%%%%%%%%%%%%%%
\subsection{Partial Sums}
\label{subsec:gf-partial-sums}
%%%%%%%%%%%%%%%%%%%%%%%%%%%%%%%%%%%%%%%%%%%%%%%%%%%%%%%%%%%%

Suppose

\[
s_n
=
\sum_{k=0}^{n}a_k.
\]

Let

\[
A(x)
=
\sum_{n\geq0}a_nx^n
\]

and

\[
S(x)
=
\sum_{n\geq0}s_nx^n.
\]

Then

\[
\boxed{
S(x)
=
\frac{A(x)}{1-x}.
}
\]

Multiplication by

\[
\frac{1}{1-x}
\]

therefore performs cumulative summation on coefficients.

%%%%%%%%%%%%%%%%%%%%%%%%%%%%%%%%%%%%%%%%%%%%%%%%%%%%%%%%%%%%
\subsection{Worked Example: Partial Sums}
\label{subsec:worked-partial-sums}
%%%%%%%%%%%%%%%%%%%%%%%%%%%%%%%%%%%%%%%%%%%%%%%%%%%%%%%%%%%%

\begin{workedexamplebox}{Partial Sums of the Constant Sequence}

Let

\[
a_n=1.
\]

Then

\[
A(x)=\frac{1}{1-x}.
\]

The partial sums are

\[
s_n=n+1.
\]

Using the partial-sum rule,

\[
S(x)
=
\frac{1}{(1-x)^2}.
\]

\begin{answerbox}
\[
\boxed{
\sum_{n=0}^{\infty}(n+1)x^n
=
\frac{1}{(1-x)^2}.
}
\]
\end{answerbox}

\end{workedexamplebox}

%%%%%%%%%%%%%%%%%%%%%%%%%%%%%%%%%%%%%%%%%%%%%%%%%%%%%%%%%%%%
\subsection{Solving Recurrences with Generating Functions}
\label{subsec:solving-recurrences-gf}
%%%%%%%%%%%%%%%%%%%%%%%%%%%%%%%%%%%%%%%%%%%%%%%%%%%%%%%%%%%%

Generating functions can transform a recurrence into an algebraic
equation.

The general procedure is:

\begin{enumerate}
    \item define the generating function;
    \item multiply the recurrence by \(x^n\);
    \item sum over all valid \(n\);
    \item rewrite shifted sums in terms of the generating function;
    \item solve algebraically for the generating function;
    \item extract coefficients.
\end{enumerate}

%%%%%%%%%%%%%%%%%%%%%%%%%%%%%%%%%%%%%%%%%%%%%%%%%%%%%%%%%%%%
\subsection{Worked Example: First-Order Recurrence}
\label{subsec:worked-gf-first-order}
%%%%%%%%%%%%%%%%%%%%%%%%%%%%%%%%%%%%%%%%%%%%%%%%%%%%%%%%%%%%

\begin{workedexamplebox}{Solving \(a_n=2a_{n-1}+1\)}

Suppose

\[
a_0=0
\]

and

\[
a_n=2a_{n-1}+1
\qquad
(n\geq1).
\]

Define

\[
A(x)
=
\sum_{n=0}^{\infty}a_nx^n.
\]

Multiply the recurrence by \(x^n\) and sum for \(n\geq1\):

\[
\sum_{n\geq1}a_nx^n
=
2\sum_{n\geq1}a_{n-1}x^n
+
\sum_{n\geq1}x^n.
\]

The left side is

\[
A(x)-a_0=A(x).
\]

The first sum on the right is

\[
2xA(x).
\]

The final sum is

\[
\frac{x}{1-x}.
\]

Therefore,

\[
A(x)
=
2xA(x)
+
\frac{x}{1-x}.
\]

Thus,

\[
A(x)(1-2x)
=
\frac{x}{1-x}.
\]

Hence,

\[
A(x)
=
\frac{x}{(1-x)(1-2x)}.
\]

Partial fractions give

\[
A(x)
=
-\frac{1}{1-x}
+
\frac{1}{1-2x}.
\]

Therefore,

\begin{answerbox}
\[
\boxed{
a_n=2^n-1.
}
\]
\end{answerbox}

\end{workedexamplebox}

%%%%%%%%%%%%%%%%%%%%%%%%%%%%%%%%%%%%%%%%%%%%%%%%%%%%%%%%%%%%
\subsection{Fibonacci Generating Function}
\label{subsec:fibonacci-generating-function}
%%%%%%%%%%%%%%%%%%%%%%%%%%%%%%%%%%%%%%%%%%%%%%%%%%%%%%%%%%%%

Let

\[
F_0=0,
\qquad
F_1=1,
\]

and

\[
F_n=F_{n-1}+F_{n-2}.
\]

Define

\[
F(x)
=
\sum_{n=0}^{\infty}F_nx^n.
\]

Using the recurrence,

\[
F(x)-x
=
xF(x)+x^2F(x).
\]

Therefore,

\[
F(x)(1-x-x^2)=x.
\]

Hence,

\[
\boxed{
F(x)
=
\frac{x}{1-x-x^2}.
}
\]

This compact rational function encodes the entire Fibonacci sequence.

%%%%%%%%%%%%%%%%%%%%%%%%%%%%%%%%%%%%%%%%%%%%%%%%%%%%%%%%%%%%
\subsection{Worked Example: Fibonacci Recurrence}
\label{subsec:worked-fibonacci-gf}
%%%%%%%%%%%%%%%%%%%%%%%%%%%%%%%%%%%%%%%%%%%%%%%%%%%%%%%%%%%%

\begin{workedexamplebox}{Recovering the Fibonacci Generating Function}

Starting from

\[
F_n=F_{n-1}+F_{n-2}
\qquad(n\geq2),
\]

multiply by \(x^n\) and sum:

\[
\sum_{n\geq2}F_nx^n
=
\sum_{n\geq2}F_{n-1}x^n
+
\sum_{n\geq2}F_{n-2}x^n.
\]

The three terms become

\[
F(x)-x,
\]

\[
xF(x),
\]

and

\[
x^2F(x).
\]

Therefore,

\[
F(x)-x
=
xF(x)+x^2F(x).
\]

\begin{answerbox}
\[
\boxed{
F(x)=\frac{x}{1-x-x^2}.
}
\]
\end{answerbox}

\end{workedexamplebox}

%%%%%%%%%%%%%%%%%%%%%%%%%%%%%%%%%%%%%%%%%%%%%%%%%%%%%%%%%%%%
\subsection{Partial Fractions}
\label{subsec:partial-fractions-gf}
%%%%%%%%%%%%%%%%%%%%%%%%%%%%%%%%%%%%%%%%%%%%%%%%%%%%%%%%%%%%

Rational generating functions can often be decomposed using partial
fractions.

Suppose

\[
A(x)
=
\frac{1}{(1-\alpha x)(1-\beta x)}.
\]

One may write

\[
A(x)
=
\frac{C}{1-\alpha x}
+
\frac{D}{1-\beta x}.
\]

Then

\[
[x^n]A(x)
=
C\alpha^n+D\beta^n.
\]

Thus factorization of the denominator reveals exponential terms in the
coefficient sequence.

%%%%%%%%%%%%%%%%%%%%%%%%%%%%%%%%%%%%%%%%%%%%%%%%%%%%%%%%%%%%
\subsection{Characteristic Roots and Generating Functions}
\label{subsec:characteristic-roots-gf}
%%%%%%%%%%%%%%%%%%%%%%%%%%%%%%%%%%%%%%%%%%%%%%%%%%%%%%%%%%%%

For constant-coefficient linear recurrences, the characteristic
polynomial and the denominator of the generating function contain the
same structural information.

For example,

\[
a_n
=
c_1a_{n-1}
+
c_2a_{n-2}
\]

has characteristic equation

\[
r^2-c_1r-c_2=0.
\]

The generating function typically has denominator

\[
\boxed{
1-c_1x-c_2x^2.
}
\]

The roots are related through reciprocal transformations.

This explains why characteristic-root methods and generating-function
methods produce equivalent formulas.

%%%%%%%%%%%%%%%%%%%%%%%%%%%%%%%%%%%%%%%%%%%%%%%%%%%%%%%%%%%%
\subsection{Repeated Factors}
\label{subsec:gf-repeated-factors}
%%%%%%%%%%%%%%%%%%%%%%%%%%%%%%%%%%%%%%%%%%%%%%%%%%%%%%%%%%%%

Repeated denominator factors correspond to polynomial factors in the
coefficients.

For example,

\[
\frac{1}{(1-rx)^2}
\]

has coefficients

\[
(n+1)r^n.
\]

More generally,

\[
\boxed{
\frac{1}{(1-rx)^k}
=
\sum_{n=0}^{\infty}
\binom{n+k-1}{k-1}
r^n x^n.
}
\]

This mirrors repeated characteristic roots in recurrence relations.

%%%%%%%%%%%%%%%%%%%%%%%%%%%%%%%%%%%%%%%%%%%%%%%%%%%%%%%%%%%%
\subsection{Generating Functions for Restricted Selections}
\label{subsec:gf-restricted-selections}
%%%%%%%%%%%%%%%%%%%%%%%%%%%%%%%%%%%%%%%%%%%%%%%%%%%%%%%%%%%%

Generating functions provide a systematic way to encode allowed
quantities.

Suppose one may choose any nonnegative number of identical objects of
one type.

Its contribution is

\[
1+x+x^2+x^3+\cdots
=
\frac{1}{1-x}.
\]

If at most two objects may be selected, the factor becomes

\[
\boxed{
1+x+x^2.
}
\]

If only an even number may be selected,

\[
\boxed{
1+x^2+x^4+\cdots
=
\frac{1}{1-x^2}.
}
\]

%%%%%%%%%%%%%%%%%%%%%%%%%%%%%%%%%%%%%%%%%%%%%%%%%%%%%%%%%%%%
\subsection{Product Rule for Selection Problems}
\label{subsec:gf-selection-product}
%%%%%%%%%%%%%%%%%%%%%%%%%%%%%%%%%%%%%%%%%%%%%%%%%%%%%%%%%%%%

Suppose several independent object types are available.

For each type, construct a generating function whose exponent records
how many objects of that type are chosen.

Multiply the factors.

Then the coefficient

\[
[x^n]
\]

counts selections containing exactly \(n\) total objects.

This provides an algebraic version of the multiplication principle.

%%%%%%%%%%%%%%%%%%%%%%%%%%%%%%%%%%%%%%%%%%%%%%%%%%%%%%%%%%%%
\subsection{Worked Example: Restricted Distribution}
\label{subsec:worked-gf-restricted-distribution}
%%%%%%%%%%%%%%%%%%%%%%%%%%%%%%%%%%%%%%%%%%%%%%%%%%%%%%%%%%%%

\begin{workedexamplebox}{Choosing Objects with Restrictions}

Suppose we choose objects of three types:

\begin{itemize}
    \item any number of type \(A\);
    \item at most two of type \(B\);
    \item an even number of type \(C\).
\end{itemize}

The generating functions are

\[
\frac{1}{1-x},
\]

\[
1+x+x^2,
\]

and

\[
\frac{1}{1-x^2}.
\]

Therefore the generating function for all valid selections is

\[
\boxed{
G(x)
=
\frac{1+x+x^2}{(1-x)(1-x^2)}.
}
\]

\begin{answerbox}
The number of valid selections using exactly \(n\) objects is

\[
\boxed{
[x^n]G(x).
}
\]
\end{answerbox}

\end{workedexamplebox}

%%%%%%%%%%%%%%%%%%%%%%%%%%%%%%%%%%%%%%%%%%%%%%%%%%%%%%%%%%%%
\subsection{Integer Solutions via Generating Functions}
\label{subsec:integer-solutions-gf}
%%%%%%%%%%%%%%%%%%%%%%%%%%%%%%%%%%%%%%%%%%%%%%%%%%%%%%%%%%%%

Consider

\[
x_1+x_2+\cdots+x_k=n
\]

with

\[
x_i\geq0.
\]

Each variable contributes

\[
1+x+x^2+\cdots
=
\frac{1}{1-x}.
\]

Therefore the total generating function is

\[
\frac{1}{(1-x)^k}.
\]

Hence,

\[
[x^n]\frac{1}{(1-x)^k}
=
\binom{n+k-1}{k-1}.
\]

Thus the stars-and-bars formula follows immediately.

%%%%%%%%%%%%%%%%%%%%%%%%%%%%%%%%%%%%%%%%%%%%%%%%%%%%%%%%%%%%
\subsection{Positive Integer Solutions via Generating Functions}
\label{subsec:positive-solutions-gf}
%%%%%%%%%%%%%%%%%%%%%%%%%%%%%%%%%%%%%%%%%%%%%%%%%%%%%%%%%%%%

If

\[
x_i\geq1,
\]

then each variable contributes

\[
x+x^2+x^3+\cdots
=
\frac{x}{1-x}.
\]

Therefore,

\[
\left(
\frac{x}{1-x}
\right)^k
=
\frac{x^k}{(1-x)^k}.
\]

The coefficient of \(x^n\) is

\[
\boxed{
\binom{n-1}{k-1}.
}
\]

This agrees with stars and bars.

%%%%%%%%%%%%%%%%%%%%%%%%%%%%%%%%%%%%%%%%%%%%%%%%%%%%%%%%%%%%
\subsection{Upper-Bounded Integer Solutions}
\label{subsec:upper-bounded-solutions-gf}
%%%%%%%%%%%%%%%%%%%%%%%%%%%%%%%%%%%%%%%%%%%%%%%%%%%%%%%%%%%%

Suppose

\[
0\leq x_i\leq m.
\]

Then each variable contributes

\[
1+x+x^2+\cdots+x^m.
\]

Therefore, solutions of

\[
x_1+\cdots+x_k=n
\]

with these bounds are counted by

\[
\boxed{
[x^n]
(1+x+\cdots+x^m)^k.
}
\]

Unlike basic stars and bars, upper bounds are incorporated directly into
the generating function.

%%%%%%%%%%%%%%%%%%%%%%%%%%%%%%%%%%%%%%%%%%%%%%%%%%%%%%%%%%%%
\subsection{Integer Partitions}
\label{subsec:integer-partitions-gf}
%%%%%%%%%%%%%%%%%%%%%%%%%%%%%%%%%%%%%%%%%%%%%%%%%%%%%%%%%%%%

An integer partition of \(n\) is a representation of \(n\) as a sum of
positive integers in which order does not matter.

For example, the partitions of \(4\) are

\[
4,
\]

\[
3+1,
\]

\[
2+2,
\]

\[
2+1+1,
\]

\[
1+1+1+1.
\]

Thus,

\[
p(4)=5.
\]

The function \(p(n)\) denotes the number of integer partitions of \(n\).

%%%%%%%%%%%%%%%%%%%%%%%%%%%%%%%%%%%%%%%%%%%%%%%%%%%%%%%%%%%%
\subsection{Euler's Partition Generating Function}
\label{subsec:euler-partition-gf}
%%%%%%%%%%%%%%%%%%%%%%%%%%%%%%%%%%%%%%%%%%%%%%%%%%%%%%%%%%%%

For each positive integer \(k\), any number of parts of size \(k\) may
appear.

Therefore the factor for parts of size \(k\) is

\[
1+x^k+x^{2k}+x^{3k}+\cdots
=
\frac{1}{1-x^k}.
\]

Multiplying over all \(k\geq1\) gives

\[
\boxed{
\sum_{n=0}^{\infty}p(n)x^n
=
\prod_{k=1}^{\infty}
\frac{1}{1-x^k}.
}
\]

This is Euler's generating function for integer partitions.

%%%%%%%%%%%%%%%%%%%%%%%%%%%%%%%%%%%%%%%%%%%%%%%%%%%%%%%%%%%%
\subsection{Partitions into Distinct Parts}
\label{subsec:partitions-distinct-parts}
%%%%%%%%%%%%%%%%%%%%%%%%%%%%%%%%%%%%%%%%%%%%%%%%%%%%%%%%%%%%

If each positive integer may appear at most once, the factor for part
size \(k\) is

\[
1+x^k.
\]

Therefore the generating function is

\[
\boxed{
\prod_{k=1}^{\infty}(1+x^k).
}
\]

The coefficient of \(x^n\) counts partitions of \(n\) into distinct
parts.

%%%%%%%%%%%%%%%%%%%%%%%%%%%%%%%%%%%%%%%%%%%%%%%%%%%%%%%%%%%%
\subsection{Partitions into Odd Parts}
\label{subsec:partitions-odd-parts}
%%%%%%%%%%%%%%%%%%%%%%%%%%%%%%%%%%%%%%%%%%%%%%%%%%%%%%%%%%%%

If only odd part sizes are allowed, the generating function becomes

\[
\boxed{
\prod_{j=1}^{\infty}
\frac{1}{1-x^{2j-1}}.
}
\]

A classical theorem of Euler states that

\[
\boxed{
\text{partitions into distinct parts}
=
\text{partitions into odd parts}.
}
\]

Equivalently,

\[
\boxed{
\prod_{k=1}^{\infty}(1+x^k)
=
\prod_{j=1}^{\infty}
\frac{1}{1-x^{2j-1}}.
}
\]

This identity has both algebraic and combinatorial interpretations.

%%%%%%%%%%%%%%%%%%%%%%%%%%%%%%%%%%%%%%%%%%%%%%%%%%%%%%%%%%%%
\subsection{Partitions with a Limited Number of Parts}
\label{subsec:partitions-limited-parts}
%%%%%%%%%%%%%%%%%%%%%%%%%%%%%%%%%%%%%%%%%%%%%%%%%%%%%%%%%%%%

Generating functions can track restrictions on both part sizes and the
number of parts.

A second variable may be introduced.

For example,

\[
\prod_{k=1}^{\infty}
\frac{1}{1-yx^k}
\]

uses:

\[
x
\]

to track total size and

\[
y
\]

to track the number of parts.

Then

\[
[x^ny^r]
\]

extracts partitions of \(n\) with exactly \(r\) parts.

This is an example of a multivariate generating function.

%%%%%%%%%%%%%%%%%%%%%%%%%%%%%%%%%%%%%%%%%%%%%%%%%%%%%%%%%%%%
\subsection{Multivariate Generating Functions}
\label{subsec:multivariate-generating-functions}
%%%%%%%%%%%%%%%%%%%%%%%%%%%%%%%%%%%%%%%%%%%%%%%%%%%%%%%%%%%%

A generating function may use multiple formal variables.

For example,

\[
\boxed{
F(x,y)
=
\sum_{m,n\geq0}
a_{m,n}x^my^n.
}
\]

The coefficient notation

\[
[x^my^n]F(x,y)
\]

extracts \(a_{m,n}\).

Multiple variables allow several statistics to be tracked
simultaneously.

%%%%%%%%%%%%%%%%%%%%%%%%%%%%%%%%%%%%%%%%%%%%%%%%%%%%%%%%%%%%
\subsection{Worked Example: Tracking Two Statistics}
\label{subsec:worked-bivariate-gf}
%%%%%%%%%%%%%%%%%%%%%%%%%%%%%%%%%%%%%%%%%%%%%%%%%%%%%%%%%%%%

\begin{workedexamplebox}{Binary Strings by Length and Weight}

For a single binary position:

\[
0
\]

contributes no \(1\), while

\[
1
\]

contributes one \(1\).

Represent this by

\[
1+y.
\]

For strings of length \(n\),

\[
(1+y)^n.
\]

The coefficient of \(y^k\) is

\[
\binom{n}{k}.
\]

To track length as well, sum over all \(n\):

\[
\sum_{n\geq0}(1+y)^nx^n.
\]

This is a geometric series.

\begin{answerbox}
\[
\boxed{
F(x,y)
=
\frac{1}{1-x(1+y)}.
}
\]

Moreover,

\[
\boxed{
[x^ny^k]F(x,y)
=
\binom{n}{k}.
}
\]
\end{answerbox}

\end{workedexamplebox}

%%%%%%%%%%%%%%%%%%%%%%%%%%%%%%%%%%%%%%%%%%%%%%%%%%%%%%%%%%%%
\subsection{Exponential Generating Functions}
\label{subsec:exponential-generating-functions}
%%%%%%%%%%%%%%%%%%%%%%%%%%%%%%%%%%%%%%%%%%%%%%%%%%%%%%%%%%%%

Ordinary generating functions are not always the most natural encoding.

For labeled combinatorial structures, exponential generating functions
are often preferable.

\begin{definitionbox}{Exponential Generating Function}
For a sequence \((a_n)_{n\geq0}\), the exponential generating function,
abbreviated EGF, is

\[
\boxed{
A(x)
=
\sum_{n=0}^{\infty}
a_n\frac{x^n}{n!}.
}
\]
\end{definitionbox}

The factorial normalization reflects the presence of labels.

%%%%%%%%%%%%%%%%%%%%%%%%%%%%%%%%%%%%%%%%%%%%%%%%%%%%%%%%%%%%
\subsection{Ordinary Versus Exponential Generating Functions}
\label{subsec:ogf-versus-egf}
%%%%%%%%%%%%%%%%%%%%%%%%%%%%%%%%%%%%%%%%%%%%%%%%%%%%%%%%%%%%

The same sequence may have both an OGF and an EGF.

For the constant sequence

\[
a_n=1,
\]

the OGF is

\[
\boxed{
\frac{1}{1-x}.
}
\]

The EGF is

\[
\boxed{
e^x.
}
\]

Indeed,

\[
e^x
=
\sum_{n=0}^{\infty}
\frac{x^n}{n!}.
\]

Thus the generating function depends not only on the sequence but also
on the chosen encoding convention.

%%%%%%%%%%%%%%%%%%%%%%%%%%%%%%%%%%%%%%%%%%%%%%%%%%%%%%%%%%%%
\subsection{EGF for Permutations}
\label{subsec:egf-permutations}
%%%%%%%%%%%%%%%%%%%%%%%%%%%%%%%%%%%%%%%%%%%%%%%%%%%%%%%%%%%%

There are

\[
n!
\]

permutations of \(n\) labeled objects.

Therefore the EGF is

\[
\sum_{n=0}^{\infty}
n!
\frac{x^n}{n!}
=
\sum_{n=0}^{\infty}x^n.
\]

Hence,

\[
\boxed{
\sum_{n=0}^{\infty}
n!\frac{x^n}{n!}
=
\frac{1}{1-x}.
}
\]

Although this looks like the ordinary generating function of the
constant sequence, here the coefficients are interpreted under EGF
normalization.

%%%%%%%%%%%%%%%%%%%%%%%%%%%%%%%%%%%%%%%%%%%%%%%%%%%%%%%%%%%%
\subsection{EGF Products}
\label{subsec:egf-products}
%%%%%%%%%%%%%%%%%%%%%%%%%%%%%%%%%%%%%%%%%%%%%%%%%%%%%%%%%%%%

Suppose

\[
A(x)
=
\sum_{n\geq0}
a_n\frac{x^n}{n!}
\]

and

\[
B(x)
=
\sum_{n\geq0}
b_n\frac{x^n}{n!}.
\]

Then

\[
A(x)B(x)
=
\sum_{n\geq0}
\left(
\sum_{k=0}^{n}
\binom{n}{k}
a_kb_{n-k}
\right)
\frac{x^n}{n!}.
\]

Thus EGF multiplication naturally introduces binomial coefficients.

This matches the combinatorial process of distributing \(n\) labels
between two labeled structures.

%%%%%%%%%%%%%%%%%%%%%%%%%%%%%%%%%%%%%%%%%%%%%%%%%%%%%%%%%%%%
\subsection{Labeled Product Principle}
\label{subsec:labeled-product-principle}
%%%%%%%%%%%%%%%%%%%%%%%%%%%%%%%%%%%%%%%%%%%%%%%%%%%%%%%%%%%%

Suppose a labeled object of total size \(n\) consists of:

\begin{itemize}
    \item one structure of size \(k\);
    \item one structure of size \(n-k\).
\end{itemize}

First choose which \(k\) labels belong to the first structure:

\[
\binom{n}{k}.
\]

Then choose the structures:

\[
a_kb_{n-k}.
\]

Thus the total is

\[
\sum_{k=0}^{n}
\binom{n}{k}
a_kb_{n-k}.
\]

This explains why EGFs are natural for labeled combinatorial classes.

%%%%%%%%%%%%%%%%%%%%%%%%%%%%%%%%%%%%%%%%%%%%%%%%%%%%%%%%%%%%
\subsection{EGF for Set Partitions}
\label{subsec:egf-set-partitions-preview}
%%%%%%%%%%%%%%%%%%%%%%%%%%%%%%%%%%%%%%%%%%%%%%%%%%%%%%%%%%%%

The Bell numbers

\[
B_n
\]

count partitions of an \(n\)-element labeled set.

Their exponential generating function is

\[
\boxed{
\sum_{n=0}^{\infty}
B_n\frac{x^n}{n!}
=
e^{e^x-1}.
}
\]

This identity is a classic example of the power of exponential
generating functions.

The derivation will be revisited in enumerative combinatorics.

%%%%%%%%%%%%%%%%%%%%%%%%%%%%%%%%%%%%%%%%%%%%%%%%%%%%%%%%%%%%
\subsection{EGF for Derangements}
\label{subsec:egf-derangements}
%%%%%%%%%%%%%%%%%%%%%%%%%%%%%%%%%%%%%%%%%%%%%%%%%%%%%%%%%%%%

Let \(D_n\) denote the number of derangements of \(n\) labeled objects.

The EGF is

\[
\boxed{
\sum_{n=0}^{\infty}
D_n\frac{x^n}{n!}
=
\frac{e^{-x}}{1-x}.
}
\]

This compact expression encodes the inclusion--exclusion formula

\[
D_n
=
n!
\sum_{k=0}^{n}\frac{(-1)^k}{k!}.
\]

%%%%%%%%%%%%%%%%%%%%%%%%%%%%%%%%%%%%%%%%%%%%%%%%%%%%%%%%%%%%
\subsection{Generating Functions and Combinatorial Classes}
\label{subsec:gf-combinatorial-classes}
%%%%%%%%%%%%%%%%%%%%%%%%%%%%%%%%%%%%%%%%%%%%%%%%%%%%%%%%%%%%

Let

\[
\mathcal{C}
\]

be a combinatorial class.

If

\[
c_n
=
|\mathcal{C}_n|,
\]

then the ordinary generating function is

\[
\boxed{
C(x)
=
\sum_{n\geq0}c_nx^n.
}
\]

Structural constructions on combinatorial classes often translate into
algebraic operations on generating functions.

This correspondence is one of the central ideas of enumerative
combinatorics.

%%%%%%%%%%%%%%%%%%%%%%%%%%%%%%%%%%%%%%%%%%%%%%%%%%%%%%%%%%%%
\subsection{Disjoint Union}
\label{subsec:gf-disjoint-union}
%%%%%%%%%%%%%%%%%%%%%%%%%%%%%%%%%%%%%%%%%%%%%%%%%%%%%%%%%%%%

Suppose

\[
\mathcal{C}
=
\mathcal{A}\sqcup\mathcal{B},
\]

where the union is disjoint.

Then

\[
c_n=a_n+b_n.
\]

Therefore,

\[
\boxed{
C(x)
=
A(x)+B(x).
}
\]

Thus disjoint union corresponds to addition.

%%%%%%%%%%%%%%%%%%%%%%%%%%%%%%%%%%%%%%%%%%%%%%%%%%%%%%%%%%%%
\subsection{Cartesian Product of Classes}
\label{subsec:gf-product-classes}
%%%%%%%%%%%%%%%%%%%%%%%%%%%%%%%%%%%%%%%%%%%%%%%%%%%%%%%%%%%%

Suppose an object in \(\mathcal{C}\) consists of an ordered pair

\[
(a,b),
\]

where

\[
a\in\mathcal{A},
\qquad
b\in\mathcal{B},
\]

and size is additive:

\[
|c|=|a|+|b|.
\]

Then

\[
c_n
=
\sum_{k=0}^{n}a_kb_{n-k}.
\]

Therefore,

\[
\boxed{
C(x)=A(x)B(x).
}
\]

Thus product of combinatorial classes corresponds to product of OGFs.

%%%%%%%%%%%%%%%%%%%%%%%%%%%%%%%%%%%%%%%%%%%%%%%%%%%%%%%%%%%%
\subsection{Sequences of Structures}
\label{subsec:gf-sequences-structures}
%%%%%%%%%%%%%%%%%%%%%%%%%%%%%%%%%%%%%%%%%%%%%%%%%%%%%%%%%%%%

Suppose \(\mathcal{A}\) contains no size-zero object.

An arbitrary finite sequence of \(\mathcal{A}\)-objects has generating
function

\[
\boxed{
1+A(x)+A(x)^2+A(x)^3+\cdots
=
\frac{1}{1-A(x)}.
}
\]

The initial \(1\) represents the empty sequence.

This is a fundamental symbolic construction.

%%%%%%%%%%%%%%%%%%%%%%%%%%%%%%%%%%%%%%%%%%%%%%%%%%%%%%%%%%%%
\subsection{Recursive Specifications}
\label{subsec:gf-recursive-specifications}
%%%%%%%%%%%%%%%%%%%%%%%%%%%%%%%%%%%%%%%%%%%%%%%%%%%%%%%%%%%%

A recursively defined combinatorial class may produce a functional
equation for its generating function.

For example, suppose

\[
\mathcal{C}
=
\{\text{empty object}\}
\cup
\left(
\text{atom}
\times
\mathcal{C}
\times
\mathcal{C}
\right).
\]

If the atom contributes size \(1\), then

\[
\boxed{
C(x)
=
1+xC(x)^2.
}
\]

This is the generating-function equation underlying Catalan structures.

%%%%%%%%%%%%%%%%%%%%%%%%%%%%%%%%%%%%%%%%%%%%%%%%%%%%%%%%%%%%
\subsection{Catalan Generating Function}
\label{subsec:catalan-generating-function}
%%%%%%%%%%%%%%%%%%%%%%%%%%%%%%%%%%%%%%%%%%%%%%%%%%%%%%%%%%%%

Recall the Catalan recurrence

\[
C_0=1
\]

and

\[
C_{n+1}
=
\sum_{k=0}^{n}C_kC_{n-k}.
\]

Let

\[
C(x)
=
\sum_{n=0}^{\infty}C_nx^n.
\]

The recurrence translates to

\[
\boxed{
C(x)
=
1+xC(x)^2.
}
\]

This is a quadratic equation in \(C(x)\).

Solving,

\[
xC(x)^2-C(x)+1=0.
\]

Using the quadratic formula,

\[
C(x)
=
\frac{1\pm\sqrt{1-4x}}{2x}.
\]

The correct branch is determined by

\[
C(0)=1.
\]

Therefore,

\[
\boxed{
C(x)
=
\frac{1-\sqrt{1-4x}}{2x}.
}
\]

%%%%%%%%%%%%%%%%%%%%%%%%%%%%%%%%%%%%%%%%%%%%%%%%%%%%%%%%%%%%
\subsection{Catalan Coefficients}
\label{subsec:catalan-coefficients}
%%%%%%%%%%%%%%%%%%%%%%%%%%%%%%%%%%%%%%%%%%%%%%%%%%%%%%%%%%%%

Expanding the Catalan generating function gives

\[
\boxed{
C_n
=
\frac{1}{n+1}\binom{2n}{n}.
}
\]

Thus

\[
1,1,2,5,14,42,132,\ldots
\]

are encoded by

\[
\frac{1-\sqrt{1-4x}}{2x}.
\]

This provides a major example of a nonlinear recurrence solved using a
generating function.

%%%%%%%%%%%%%%%%%%%%%%%%%%%%%%%%%%%%%%%%%%%%%%%%%%%%%%%%%%%%
\subsection{Worked Example: Catalan Functional Equation}
\label{subsec:worked-catalan-gf}
%%%%%%%%%%%%%%%%%%%%%%%%%%%%%%%%%%%%%%%%%%%%%%%%%%%%%%%%%%%%

\begin{workedexamplebox}{From Recurrence to Functional Equation}

Suppose

\[
C_0=1
\]

and

\[
C_{n+1}
=
\sum_{k=0}^{n}C_kC_{n-k}.
\]

Multiply by \(x^{n+1}\) and sum for \(n\geq0\):

\[
\sum_{n\geq0}C_{n+1}x^{n+1}
=
x
\sum_{n\geq0}
\left(
\sum_{k=0}^{n}C_kC_{n-k}
\right)x^n.
\]

The left side is

\[
C(x)-1.
\]

The convolution on the right is the coefficient sequence of

\[
C(x)^2.
\]

Therefore,

\[
C(x)-1=xC(x)^2.
\]

\begin{answerbox}
\[
\boxed{
C(x)=1+xC(x)^2.
}
\]
\end{answerbox}

\end{workedexamplebox}

%%%%%%%%%%%%%%%%%%%%%%%%%%%%%%%%%%%%%%%%%%%%%%%%%%%%%%%%%%%%
\subsection{Rational Generating Functions}
\label{subsec:rational-generating-functions}
%%%%%%%%%%%%%%%%%%%%%%%%%%%%%%%%%%%%%%%%%%%%%%%%%%%%%%%%%%%%

A generating function is rational if it can be written

\[
\boxed{
A(x)=\frac{P(x)}{Q(x)},
}
\]

where \(P(x)\) and \(Q(x)\) are polynomials.

Sequences satisfying linear recurrences with constant coefficients
typically have rational generating functions.

Conversely, coefficients of rational generating functions eventually
satisfy linear recurrences with constant coefficients.

Thus rational generating functions and constant-coefficient linear
recurrences are closely connected.

%%%%%%%%%%%%%%%%%%%%%%%%%%%%%%%%%%%%%%%%%%%%%%%%%%%%%%%%%%%%
\subsection{Algebraic Generating Functions}
\label{subsec:algebraic-generating-functions}
%%%%%%%%%%%%%%%%%%%%%%%%%%%%%%%%%%%%%%%%%%%%%%%%%%%%%%%%%%%%

A generating function \(A(x)\) is algebraic if it satisfies a
polynomial equation

\[
\boxed{
P(x,A(x))=0
}
\]

for some nonzero polynomial \(P\).

The Catalan generating function is algebraic because it satisfies

\[
xC(x)^2-C(x)+1=0.
\]

Every rational generating function is algebraic, but not every
algebraic generating function is rational.

%%%%%%%%%%%%%%%%%%%%%%%%%%%%%%%%%%%%%%%%%%%%%%%%%%%%%%%%%%%%
\subsection{Transcendental Generating Functions}
\label{subsec:transcendental-generating-functions}
%%%%%%%%%%%%%%%%%%%%%%%%%%%%%%%%%%%%%%%%%%%%%%%%%%%%%%%%%%%%

Some generating functions are not algebraic.

For example,

\[
e^x
\]

is transcendental over the rational function field.

Exponential generating functions frequently involve functions such as

\[
e^x,
\qquad
\log(1-x),
\qquad
\sin x,
\qquad
\cos x.
\]

These arise naturally when labeled combinatorial structures are
assembled.

%%%%%%%%%%%%%%%%%%%%%%%%%%%%%%%%%%%%%%%%%%%%%%%%%%%%%%%%%%%%
\subsection{Generating Functions and Probability}
\label{subsec:generating-functions-probability-preview}
%%%%%%%%%%%%%%%%%%%%%%%%%%%%%%%%%%%%%%%%%%%%%%%%%%%%%%%%%%%%

Probability theory uses several related generating transforms.

For a nonnegative integer-valued random variable \(X\), the probability
generating function is

\[
\boxed{
G_X(s)
=
\mathbb{E}[s^X].
}
\]

If

\[
\mathbb{P}(X=n)=p_n,
\]

then

\[
G_X(s)
=
\sum_{n=0}^{\infty}p_ns^n.
\]

Thus a probability generating function is an ordinary generating
function whose coefficients are probabilities.

This topic is developed more fully in Chapter~\ref{chap:probability}.

%%%%%%%%%%%%%%%%%%%%%%%%%%%%%%%%%%%%%%%%%%%%%%%%%%%%%%%%%%%%
\subsection{Moment Generating Functions}
\label{subsec:moment-generating-functions-preview}
%%%%%%%%%%%%%%%%%%%%%%%%%%%%%%%%%%%%%%%%%%%%%%%%%%%%%%%%%%%%

Probability also uses the moment generating function

\[
\boxed{
M_X(t)
=
\mathbb{E}[e^{tX}].
}
\]

This is not the same object as an ordinary combinatorial generating
function, although the underlying philosophy is similar:

\[
\boxed{
\text{encode a sequence or distribution into a function}.
}
\]

Different generating transforms are chosen according to the structure
being studied.

%%%%%%%%%%%%%%%%%%%%%%%%%%%%%%%%%%%%%%%%%%%%%%%%%%%%%%%%%%%%
\subsection{Generating Functions and Linear Algebra}
\label{subsec:generating-functions-linear-algebra}
%%%%%%%%%%%%%%%%%%%%%%%%%%%%%%%%%%%%%%%%%%%%%%%%%%%%%%%%%%%%

Generating functions can also encode matrix sequences.

Suppose \(A\) is a matrix.

Then

\[
I+xA+x^2A^2+x^3A^3+\cdots
\]

formally satisfies

\[
\boxed{
(I-xA)^{-1}
=
\sum_{n=0}^{\infty}x^nA^n.
}
\]

This identity connects matrix powers, walks in graphs, linear
recurrences, and spectral methods.

%%%%%%%%%%%%%%%%%%%%%%%%%%%%%%%%%%%%%%%%%%%%%%%%%%%%%%%%%%%%
\subsection{Graph Walks and Generating Functions}
\label{subsec:graph-walks-gf-preview}
%%%%%%%%%%%%%%%%%%%%%%%%%%%%%%%%%%%%%%%%%%%%%%%%%%%%%%%%%%%%

If \(A\) is the adjacency matrix of a finite graph, then

\[
(A^n)_{ij}
\]

counts walks of length \(n\) from vertex \(i\) to vertex \(j\).

Therefore,

\[
\boxed{
\sum_{n=0}^{\infty}
(A^n)_{ij}x^n
}
\]

generates the number of walks of each length.

Matrix generating functions thus provide a natural bridge to graph
theory.

%%%%%%%%%%%%%%%%%%%%%%%%%%%%%%%%%%%%%%%%%%%%%%%%%%%%%%%%%%%%
\subsection{Coefficient Extraction Strategies}
\label{subsec:coefficient-extraction-strategies}
%%%%%%%%%%%%%%%%%%%%%%%%%%%%%%%%%%%%%%%%%%%%%%%%%%%%%%%%%%%%

Given a generating function \(A(x)\), common coefficient-extraction
methods include:

\begin{itemize}
    \item recognizing a standard generating function;
    \item geometric-series expansion;
    \item binomial-series expansion;
    \item partial fractions;
    \item differentiation;
    \item convolution;
    \item recurrence extraction;
    \item algebraic series expansion.
\end{itemize}

The appropriate method depends on the algebraic form of \(A(x)\).

%%%%%%%%%%%%%%%%%%%%%%%%%%%%%%%%%%%%%%%%%%%%%%%%%%%%%%%%%%%%
\subsection{Worked Example: Partial Fraction Extraction}
\label{subsec:worked-partial-fraction-extraction}
%%%%%%%%%%%%%%%%%%%%%%%%%%%%%%%%%%%%%%%%%%%%%%%%%%%%%%%%%%%%

\begin{workedexamplebox}{Extracting a Coefficient}

Consider

\[
A(x)
=
\frac{1}{(1-x)(1-2x)}.
\]

Write

\[
\frac{1}{(1-x)(1-2x)}
=
-\frac{1}{1-x}
+
\frac{2}{1-2x}.
\]

Using geometric expansions,

\[
-\frac{1}{1-x}
=
-\sum_{n\geq0}x^n,
\]

while

\[
\frac{2}{1-2x}
=
2\sum_{n\geq0}2^nx^n.
\]

Therefore,

\begin{answerbox}
\[
\boxed{
[x^n]A(x)
=
2^{n+1}-1.
}
\]
\end{answerbox}

\end{workedexamplebox}

%%%%%%%%%%%%%%%%%%%%%%%%%%%%%%%%%%%%%%%%%%%%%%%%%%%%%%%%%%%%
\subsection{Generating Functions for Periodic Sequences}
\label{subsec:gf-periodic-sequences}
%%%%%%%%%%%%%%%%%%%%%%%%%%%%%%%%%%%%%%%%%%%%%%%%%%%%%%%%%%%%

Suppose a sequence repeats with period \(m\):

\[
a_{n+m}=a_n.
\]

Then

\[
A(x)
=
(a_0+a_1x+\cdots+a_{m-1}x^{m-1})
(1+x^m+x^{2m}+\cdots).
\]

Therefore,

\[
\boxed{
A(x)
=
\frac{
a_0+a_1x+\cdots+a_{m-1}x^{m-1}
}{
1-x^m
}.
}
\]

Thus periodic sequences have rational generating functions.

%%%%%%%%%%%%%%%%%%%%%%%%%%%%%%%%%%%%%%%%%%%%%%%%%%%%%%%%%%%%
\subsection{Worked Example: Alternating Sequence}
\label{subsec:worked-alternating-gf}
%%%%%%%%%%%%%%%%%%%%%%%%%%%%%%%%%%%%%%%%%%%%%%%%%%%%%%%%%%%%

\begin{workedexamplebox}{Generating Function for \(1,-1,1,-1,\ldots\)}

The sequence is

\[
a_n=(-1)^n.
\]

Its generating function is

\[
1-x+x^2-x^3+\cdots.
\]

This is geometric with ratio \(-x\).

\begin{answerbox}
\[
\boxed{
A(x)=\frac{1}{1+x}.
}
\]
\end{answerbox}

\end{workedexamplebox}

%%%%%%%%%%%%%%%%%%%%%%%%%%%%%%%%%%%%%%%%%%%%%%%%%%%%%%%%%%%%
\subsection{Finite Generating Functions}
\label{subsec:finite-generating-functions}
%%%%%%%%%%%%%%%%%%%%%%%%%%%%%%%%%%%%%%%%%%%%%%%%%%%%%%%%%%%%

Generating functions need not be infinite.

For a finite sequence

\[
a_0,a_1,\ldots,a_m,
\]

one may use the polynomial

\[
\boxed{
A(x)
=
a_0+a_1x+\cdots+a_mx^m.
}
\]

Finite generating polynomials appear naturally in:

\begin{itemize}
    \item restricted selections;
    \item graph invariants;
    \item coding theory;
    \item probability distributions;
    \item rank-generating functions.
\end{itemize}

%%%%%%%%%%%%%%%%%%%%%%%%%%%%%%%%%%%%%%%%%%%%%%%%%%%%%%%%%%%%
\subsection{Subset Generating Polynomial}
\label{subsec:subset-generating-polynomial}
%%%%%%%%%%%%%%%%%%%%%%%%%%%%%%%%%%%%%%%%%%%%%%%%%%%%%%%%%%%%

Let \(S\) be an \(n\)-element set.

Track subsets according to size.

Each element contributes either:

\[
1
\]

if excluded or

\[
x
\]

if included.

Thus the generating polynomial is

\[
\boxed{
(1+x)^n.
}
\]

The coefficient of \(x^k\) is

\[
\binom{n}{k}.
\]

Therefore,

\[
\boxed{
(1+x)^n
=
\sum_{k=0}^{n}
\binom{n}{k}x^k.
}
\]

The binomial theorem is therefore itself a generating-function identity.

%%%%%%%%%%%%%%%%%%%%%%%%%%%%%%%%%%%%%%%%%%%%%%%%%%%%%%%%%%%%
\subsection{Weighted Generating Functions}
\label{subsec:weighted-generating-functions}
%%%%%%%%%%%%%%%%%%%%%%%%%%%%%%%%%%%%%%%%%%%%%%%%%%%%%%%%%%%%

Generating functions may encode weights rather than simply counts.

Suppose each object \(c\) has size \(|c|\) and weight \(w(c)\).

Then define

\[
\boxed{
F(x)
=
\sum_{c}
w(c)x^{|c|}.
}
\]

The coefficient

\[
[x^n]F(x)
\]

is the total weight of all size-\(n\) objects.

Ordinary counting is recovered when

\[
w(c)=1.
\]

%%%%%%%%%%%%%%%%%%%%%%%%%%%%%%%%%%%%%%%%%%%%%%%%%%%%%%%%%%%%
\subsection{Sign-Weighted Generating Functions}
\label{subsec:sign-weighted-generating-functions}
%%%%%%%%%%%%%%%%%%%%%%%%%%%%%%%%%%%%%%%%%%%%%%%%%%%%%%%%%%%%

Weights may be negative.

For example,

\[
\sum_{k=0}^{n}
(-1)^k
\binom{n}{k}
x^k
=
(1-x)^n.
\]

The sign

\[
(-1)^k
\]

records the parity of the selected subset size.

Such signed generating functions appear in inclusion--exclusion,
determinants, Möbius inversion, topology, and algebraic combinatorics.

%%%%%%%%%%%%%%%%%%%%%%%%%%%%%%%%%%%%%%%%%%%%%%%%%%%%%%%%%%%%
\subsection{Generating Functions and Asymptotics}
\label{subsec:gf-asymptotics-preview}
%%%%%%%%%%%%%%%%%%%%%%%%%%%%%%%%%%%%%%%%%%%%%%%%%%%%%%%%%%%%

The analytic behavior of a generating function can reveal the growth of
its coefficients.

For example,

\[
\frac{1}{1-rx}
\]

has a singularity at

\[
x=\frac1r.
\]

Its coefficients are

\[
r^n.
\]

More generally, singularities nearest the origin often control
asymptotic coefficient growth.

This viewpoint leads to analytic combinatorics.

%%%%%%%%%%%%%%%%%%%%%%%%%%%%%%%%%%%%%%%%%%%%%%%%%%%%%%%%%%%%
\subsection{Dominant Singularities}
\label{subsec:dominant-singularities-preview}
%%%%%%%%%%%%%%%%%%%%%%%%%%%%%%%%%%%%%%%%%%%%%%%%%%%%%%%%%%%%

A singularity of smallest modulus is often called a dominant
singularity.

If a generating function behaves roughly like

\[
\frac{C}{1-x/\rho}
\]

near a dominant singularity \(\rho\), then coefficients often behave
like

\[
\boxed{
C\rho^{-n}.
}
\]

The Greek letter

\[
\rho
\]

is called ``rho.''

A rigorous treatment requires complex analysis and belongs to more
advanced enumerative methods.

%%%%%%%%%%%%%%%%%%%%%%%%%%%%%%%%%%%%%%%%%%%%%%%%%%%%%%%%%%%%
\subsection{Why Generating Functions Are Powerful}
\label{subsec:why-gf-powerful}
%%%%%%%%%%%%%%%%%%%%%%%%%%%%%%%%%%%%%%%%%%%%%%%%%%%%%%%%%%%%

Generating functions convert combinatorial information into algebra.

The basic dictionary is

\[
\boxed{
\begin{array}{c|c}
\text{Combinatorial Operation}
&
\text{Generating-Function Operation}\\
\hline
\text{disjoint choice}
&
\text{addition}\\
\text{ordered combination}
&
\text{multiplication}\\
\text{size shift}
&
\text{multiplication by }x\\
\text{sequence formation}
&
\dfrac{1}{1-A(x)}\\
\text{recurrence}
&
\text{functional equation}\\
\text{cumulative sum}
&
\dfrac{A(x)}{1-x}
\end{array}
}
\]

The strength of the method comes from this translation between
structure and algebra.

%%%%%%%%%%%%%%%%%%%%%%%%%%%%%%%%%%%%%%%%%%%%%%%%%%%%%%%%%%%%
\subsection{Common Errors}
\label{subsec:generating-functions-common-errors}
%%%%%%%%%%%%%%%%%%%%%%%%%%%%%%%%%%%%%%%%%%%%%%%%%%%%%%%%%%%%

\subsubsection{Forgetting the Initial Term}

If

\[
A(x)
=
\sum_{n\geq0}a_nx^n,
\]

then

\[
\sum_{n\geq1}a_nx^n
=
A(x)-a_0.
\]

Ignoring this correction is one of the most common recurrence-solving
errors.

%%%%%%%%%%%%%%%%%%%%%%%%%%%%%%%%%%%%%%%%%%%%%%%%%%%%%%%%%%%%
\subsubsection{Incorrect Index Shifts}

For example,

\[
\sum_{n\geq1}a_{n-1}x^n
=
xA(x),
\]

not simply \(A(x)\).

Always rewrite the exponent and sequence index carefully.

%%%%%%%%%%%%%%%%%%%%%%%%%%%%%%%%%%%%%%%%%%%%%%%%%%%%%%%%%%%%
\subsubsection{Confusing the Sequence with Its Generating Function}

The sequence is

\[
a_0,a_1,a_2,\ldots.
\]

The generating function is

\[
A(x)
=
\sum a_nx^n.
\]

They encode the same information but are different mathematical
objects.

%%%%%%%%%%%%%%%%%%%%%%%%%%%%%%%%%%%%%%%%%%%%%%%%%%%%%%%%%%%%
\subsubsection{Confusing OGF and EGF Coefficients}

For an OGF,

\[
A(x)
=
\sum a_nx^n.
\]

For an EGF,

\[
A(x)
=
\sum a_n\frac{x^n}{n!}.
\]

Coefficient interpretation therefore differs by a factorial.

%%%%%%%%%%%%%%%%%%%%%%%%%%%%%%%%%%%%%%%%%%%%%%%%%%%%%%%%%%%%
\subsubsection{Treating Formal Power Series Only as Analytic Functions}

In combinatorics, convergence may be irrelevant.

Formal power series can be manipulated coefficientwise.

Analytic convergence becomes important only when analytic methods are
introduced.

%%%%%%%%%%%%%%%%%%%%%%%%%%%%%%%%%%%%%%%%%%%%%%%%%%%%%%%%%%%%
\subsubsection{Using Product Without Checking the Size Rule}

The ordinary product corresponds to combining structures when sizes
add.

If size behaves differently, the product interpretation may not apply.

%%%%%%%%%%%%%%%%%%%%%%%%%%%%%%%%%%%%%%%%%%%%%%%%%%%%%%%%%%%%
\subsubsection{Ignoring Restrictions in Selection Factors}

If only \(0,1,\) or \(2\) objects of one type may be selected, the
correct factor is

\[
1+x+x^2,
\]

not

\[
\frac{1}{1-x}.
\]

%%%%%%%%%%%%%%%%%%%%%%%%%%%%%%%%%%%%%%%%%%%%%%%%%%%%%%%%%%%%
\subsubsection{Choosing the Wrong Branch of an Algebraic Solution}

The Catalan functional equation yields two algebraic solutions.

Only one has the correct formal power-series behavior with

\[
C(0)=1.
\]

Initial coefficient conditions determine the correct branch.

%%%%%%%%%%%%%%%%%%%%%%%%%%%%%%%%%%%%%%%%%%%%%%%%%%%%%%%%%%%%
\subsubsection{Assuming Every Generating Function Is Rational}

Some generating functions are rational.

Others are algebraic or transcendental.

The algebraic type reflects deeper structure in the sequence.

%%%%%%%%%%%%%%%%%%%%%%%%%%%%%%%%%%%%%%%%%%%%%%%%%%%%%%%%%%%%
\subsection{Applications}
\label{subsec:generating-functions-applications}
%%%%%%%%%%%%%%%%%%%%%%%%%%%%%%%%%%%%%%%%%%%%%%%%%%%%%%%%%%%%

\begin{applicationbox}

\textbf{Enumerative Combinatorics.}
Generating functions encode counting sequences and translate structural
decompositions into algebraic equations.

\medskip

\textbf{Recurrence Relations.}
Linear and nonlinear recurrences can often be solved by converting them
into functional equations.

\medskip

\textbf{Integer Partitions.}
Infinite products encode allowable part sizes and multiplicities.

\medskip

\textbf{Probability.}
Probability generating functions encode discrete probability
distributions.

\medskip

\textbf{Graph Theory.}
Matrix generating functions encode numbers of walks of each length.

\medskip

\textbf{Computer Science.}
Generating functions appear in algorithm analysis, formal languages,
recursive structures, and average-case analysis.

\medskip

\textbf{Number Theory.}
Generating series encode divisor functions, partitions, modular forms,
and arithmetic sequences.

\medskip

\textbf{Algebra.}
Hilbert series, Poincaré series, and graded dimensions are generating
functions associated with algebraic structures.

\medskip

\textbf{Statistical Mechanics.}
Partition functions are weighted generating objects encoding physical
states.

\end{applicationbox}

%%%%%%%%%%%%%%%%%%%%%%%%%%%%%%%%%%%%%%%%%%%%%%%%%%%%%%%%%%%%
\subsection{Exercises}
\label{subsec:generating-functions-exercises}
%%%%%%%%%%%%%%%%%%%%%%%%%%%%%%%%%%%%%%%%%%%%%%%%%%%%%%%%%%%%

\begin{exercise}[1]
Find the ordinary generating function of

\[
1,1,1,1,\ldots.
\]
\end{exercise}

\begin{exercise}[1]
Find the ordinary generating function of

\[
1,2,4,8,16,\ldots.
\]
\end{exercise}

\begin{exercise}[1]
Find the ordinary generating function of

\[
1,-1,1,-1,\ldots.
\]
\end{exercise}

\begin{exercise}[1]
Find

\[
[x^5]\frac{1}{1-x}.
\]
\end{exercise}

\begin{exercise}[1]
Find

\[
[x^4]\frac{1}{1-3x}.
\]
\end{exercise}

\begin{exercise}[1]
Find the generating function for the sequence

\[
0,1,2,3,4,\ldots.
\]
\end{exercise}

\begin{exercise}[1]
Show that

\[
\sum_{n\geq0}(n+1)x^n
=
\frac{1}{(1-x)^2}.
\]
\end{exercise}

\begin{exercise}[1]
Find the coefficient of \(x^7\) in

\[
(1+x)^{10}.
\]
\end{exercise}

\begin{exercise}[1]
Find the coefficient of \(x^8\) in

\[
\frac{1}{(1-x)^3}.
\]
\end{exercise}

\begin{exercise}[1]
Explain the meaning of

\[
[x^n]A(x).
\]
\end{exercise}

\begin{exercise}[2]
Derive

\[
\sum_{n\geq0}nx^n
=
\frac{x}{(1-x)^2}.
\]
\end{exercise}

\begin{exercise}[2]
Derive

\[
\sum_{n\geq0}n^2x^n
=
\frac{x(1+x)}{(1-x)^3}.
\]
\end{exercise}

\begin{exercise}[2]
Find the generating function for

\[
a_n=5^n.
\]
\end{exercise}

\begin{exercise}[2]
Find the generating function for the partial sums of

\[
1,1,1,\ldots.
\]
\end{exercise}

\begin{exercise}[2]
Solve using generating functions:

\[
a_n=3a_{n-1}+1,
\qquad
a_0=0.
\]
\end{exercise}

\begin{exercise}[2]
Derive the Fibonacci generating function

\[
F(x)=\frac{x}{1-x-x^2}.
\]
\end{exercise}

\begin{exercise}[2]
Use partial fractions to extract coefficients from

\[
\frac{1}{(1-x)(1-3x)}.
\]
\end{exercise}

\begin{exercise}[2]
Find the coefficient of \(x^{10}\) in

\[
\frac{1}{(1-x)^4}.
\]
\end{exercise}

\begin{exercise}[2]
Use generating functions to count nonnegative integer solutions of

\[
x_1+x_2+x_3=12.
\]
\end{exercise}

\begin{exercise}[2]
Use generating functions to count positive integer solutions of

\[
x_1+x_2+x_3+x_4=15.
\]
\end{exercise}

\begin{exercise}[2]
Write a generating function for solutions of

\[
x_1+x_2+x_3=n
\]

with

\[
0\leq x_i\leq4.
\]
\end{exercise}

\begin{exercise}[2]
Write a generating function for partitions using only odd parts.
\end{exercise}

\begin{exercise}[2]
Write a generating function for partitions into distinct parts.
\end{exercise}

\begin{exercise}[2]
Find the EGF of the constant sequence

\[
a_n=1.
\]
\end{exercise}

\begin{exercise}[2]
Find the EGF of

\[
a_n=n!.
\]
\end{exercise}

\begin{exercise}[3]
Solve using generating functions:

\[
a_n=2a_{n-1}+3,
\qquad
a_0=1.
\]
\end{exercise}

\begin{exercise}[3]
Solve using generating functions:

\[
a_n=5a_{n-1}-6a_{n-2},
\qquad
a_0=1,
\qquad
a_1=4.
\]
\end{exercise}

\begin{exercise}[3]
Show that a sequence satisfying a linear recurrence with constant
coefficients has a rational generating function.
\end{exercise}

\begin{exercise}[3]
Show that

\[
\frac{1}{(1-rx)^2}
=
\sum_{n\geq0}(n+1)r^nx^n.
\]
\end{exercise}

\begin{exercise}[3]
Prove using generating functions that

\[
\sum_{k=0}^{n}\binom{n}{k}=2^n.
\]
\end{exercise}

\begin{exercise}[3]
Find a bivariate generating function for binary strings where \(x\)
tracks length and \(y\) tracks the number of \(1\)'s.
\end{exercise}

\begin{exercise}[3]
Show that

\[
[x^ny^k]
\frac{1}{1-x(1+y)}
=
\binom{n}{k}.
\]
\end{exercise}

\begin{exercise}[3]
Derive the functional equation

\[
C(x)=1+xC(x)^2
\]

from the Catalan recurrence.
\end{exercise}

\begin{exercise}[3]
Solve the quadratic equation for \(C(x)\) and explain why

\[
C(x)
=
\frac{1-\sqrt{1-4x}}{2x}
\]

is the correct formal power-series branch.
\end{exercise}

\begin{exercise}[3]
Use the Catalan generating function to derive the first six Catalan
numbers.
\end{exercise}

\begin{exercise}[3]
Show that

\[
\prod_{k=1}^{\infty}
\frac{1}{1-x^k}
\]

is the generating function for integer partitions.
\end{exercise}

\begin{exercise}[3]
Explain combinatorially why

\[
\prod_{k=1}^{\infty}
(1+x^k)
\]

counts partitions into distinct parts.
\end{exercise}

\begin{exercise}[3]
Derive the EGF product formula

\[
c_n
=
\sum_{k=0}^{n}
\binom{n}{k}a_kb_{n-k}.
\]
\end{exercise}

\begin{exercise}[4]
Study the Bell-number exponential generating function

\[
e^{e^x-1}.
\]

Explain its combinatorial meaning.
\end{exercise}

\begin{exercise}[4]
Study the derangement EGF

\[
\frac{e^{-x}}{1-x}
\]

and recover the inclusion--exclusion formula for \(D_n\).
\end{exercise}

\begin{exercise}[4]
Investigate Euler's theorem equating partitions into distinct parts with
partitions into odd parts.
\end{exercise}

\begin{exercise}[4]
Study multivariate generating functions for integer partitions with a
fixed number of parts.
\end{exercise}

\begin{exercise}[4]
Investigate rational generating functions and prove their relationship
with constant-coefficient linear recurrences.
\end{exercise}

\begin{exercise}[4]
Investigate algebraic generating functions and show that the Catalan
generating function is algebraic but not rational.
\end{exercise}

\begin{exercise}[4]
Study the transfer-matrix method for counting walks in finite graphs.
\end{exercise}

\begin{exercise}[4]
Investigate probability generating functions for discrete random
variables.
\end{exercise}

\begin{exercise}[4]
Study singularity analysis and explain how dominant singularities can
determine coefficient growth.
\end{exercise}

\begin{exercise}[4]
Investigate the symbolic method of analytic combinatorics and the
translation of combinatorial class constructions into generating
functions.
\end{exercise}

%%%%%%%%%%%%%%%%%%%%%%%%%%%%%%%%%%%%%%%%%%%%%%%%%%%%%%%%%%%%
\subsection{Conceptual Summary}
\label{subsec:generating-functions-summary}
%%%%%%%%%%%%%%%%%%%%%%%%%%%%%%%%%%%%%%%%%%%%%%%%%%%%%%%%%%%%

An ordinary generating function encodes a sequence as

\[
\boxed{
A(x)
=
\sum_{n\geq0}a_nx^n.
}
\]

Coefficient extraction is written

\[
\boxed{
[x^n]A(x)=a_n.
}
\]

The geometric identity

\[
\boxed{
\frac{1}{1-x}
=
\sum_{n\geq0}x^n
}
\]

is the basic generating-function building block.

Multiplication creates convolution:

\[
\boxed{
[x^n](A(x)B(x))
=
\sum_{k=0}^{n}a_kb_{n-k}.
}
\]

Differentiation introduces index factors:

\[
\boxed{
xA'(x)
=
\sum_{n\geq0}na_nx^n.
}
\]

Partial sums satisfy

\[
\boxed{
S(x)
=
\frac{A(x)}{1-x}.
}
\]

Constant-coefficient linear recurrences usually produce rational
generating functions.

The Fibonacci sequence has generating function

\[
\boxed{
F(x)
=
\frac{x}{1-x-x^2}.
}
\]

Restricted selection problems are encoded by choosing appropriate
factors.

For example,

\[
1+x+x^2
\]

represents choosing zero, one, or two objects of a given type.

Integer partitions satisfy

\[
\boxed{
\sum_{n\geq0}p(n)x^n
=
\prod_{k\geq1}\frac{1}{1-x^k}.
}
\]

Exponential generating functions use

\[
\boxed{
A(x)
=
\sum_{n\geq0}
a_n\frac{x^n}{n!}
}
\]

and are especially natural for labeled structures.

Recursive combinatorial specifications translate into functional
equations.

For Catalan structures,

\[
\boxed{
C(x)
=
1+xC(x)^2.
}
\]

Thus generating functions form a bridge between

\[
\boxed{
\text{sequences}
+
\text{recurrences}
+
\text{counting}
+
\text{algebra}
+
\text{combinatorial structure}.
}
\]

%%%%%%%%%%%%%%%%%%%%%%%%%%%%%%%%%%%%%%%%%%%%%%%%%%%%%%%%%%%%
\subsection{Section Connections}
\label{subsec:generating-functions-connections}
%%%%%%%%%%%%%%%%%%%%%%%%%%%%%%%%%%%%%%%%%%%%%%%%%%%%%%%%%%%%

\chapterconnection{
Generating Functions extend the counting methods of
Section~\ref{sec:basic-counting} and provide an algebraic method for
solving the recurrence relations of
Section~\ref{sec:recurrence-relations}. The next section, Enumerative
Combinatorics, uses generating functions systematically for classical
families such as permutations, partitions, Stirling numbers, Bell
numbers, Catalan objects, and lattice paths. Algebraic Combinatorics
uses generating functions together with group actions, symmetric
functions, and representation-theoretic structures. Graph Theory uses
matrix and walk-generating functions. Probability uses probability and
moment generating functions. Number Theory uses generating series for
partitions and arithmetic functions, while Analysis provides the
analytic tools needed for advanced coefficient asymptotics.
}

%%%%%%%%%%%%%%%%%%%%%%%%%%%%%%%%%%%%%%%%%%%%%%%%%%%%%%%%%%%%
% SECTION SOURCE TRACKING
%%%%%%%%%%%%%%%%%%%%%%%%%%%%%%%%%%%%%%%%%%%%%%%%%%%%%%%%%%%%

%------------------------------------------------------------
% 6.4 Generating Functions
%------------------------------------------------------------
%
% Source basis:
%   - Standard enumerative combinatorics.
%   - Standard discrete mathematics.
%   - Standard recurrence theory.
%   - Standard formal power-series methods.
%
% Used for:
%   - ordinary generating functions;
%   - coefficient extraction;
%   - formal power series;
%   - geometric generating functions;
%   - sequence shifts;
%   - addition and scalar multiplication;
%   - products and convolution;
%   - negative binomial series;
%   - differentiation and integration;
%   - generating functions for n and n^2;
%   - partial sums;
%   - solving recurrence relations;
%   - Fibonacci generating functions;
%   - partial fractions;
%   - repeated denominator factors;
%   - restricted selections;
%   - integer-solution counting;
%   - integer partitions;
%   - partitions into distinct and odd parts;
%   - multivariate generating functions;
%   - exponential generating functions;
%   - labeled products;
%   - Bell-number and derangement previews;
%   - combinatorial classes;
%   - symbolic union and product constructions;
%   - sequence constructions;
%   - recursive specifications;
%   - Catalan generating functions;
%   - rational, algebraic, and transcendental generating functions;
%   - probability generating functions;
%   - matrix and graph-walk generating functions;
%   - coefficient extraction;
%   - periodic sequences;
%   - finite generating polynomials;
%   - weighted generating functions;
%   - asymptotic and singularity previews;
%   - applications and exercises.
%
% Notes:
%   - Generating functions are treated primarily as formal power series
%     in this section.
%   - Analytic convergence and singularity analysis belong to more
%     advanced analytic combinatorics.
%   - Integer partitions and classical enumerative sequences are
%     developed more fully in Section 6.5.
%   - Exponential generating functions become especially important in
%     labeled enumeration.
%
%%%%%%%%%%%%%%%%%%%%%%%%%%%%%%%%%%%%%%%%%%%%%%%%%%%%%%%%%%%%